\subsection{Sprint 1 --- Banco de dados e tasks do front-end}

Com as entregas dos mockups e as definições de escopo estabelecidas na Sprint 0,
demos seguimento à planning, onde foram definidos os Squads e as tarefas a serem
implementadas por cada grupo. No meu grupo foram sorteados Maira (\ac{ages} I) e
Vinícios (\ac{ages} III). A apreensão inicial na divisão das tarefas deixou todos nervosos,
mas, como de hábito, criei um grupo, chamei os colegas e compartilhei as tarefas para
conversarmos. Chegamos à conclusão de que eu e a Maira ficaríamos responsáveis pela
entrega das telas de \textbf{cadastro de aluno} no front-end.

Com a ajuda do Hartur (\ac{ages} IV), criei a branch \textit{feat-student-cadastro}. No
domingo da mesma semana --- que coincidiu com o domingo de Páscoa --- o Hartur
entrou em call e demonstrou como trabalhar com React no front-end, ajudando-nos a
iniciar as tasks. Ele criou um gerenciador de \textit{steps} para controlar e validar as
interações com as duas telas, e desenvolveu a tela do step 01/04. Fiquei responsável
pelo step 02/04, que comecei a desenvolver ainda no domingo e terminei na
segunda-feira antes da aula. Dei o push para averiguação do meu responsável de \ac{ages},
que verificou e me apontou pontos a melhorar. Fiz as alterações durante a aula. Com
isso, a Maira começou o step 03/04 e, para não atrasar a entrega, iniciei o step 04/04,
que ficou pronto na quinta-feira da mesma semana.

A princípio, tudo estava bem. Porém, no domingo, dia 12/04, durante algumas
integrações, o Guilherme (\ac{ages} IV) identificou problemas de segurança na aplicação.
Acompanhei a explicação onde ele demonstrou o que estava comprometendo o sistema
e o que fazer para corrigi-lo. Foi uma experiência muito enriquecedora, e pretendo
aplicar essas orientações para elevar a qualidade do meu código.

Na semana seguinte, surgiu uma demanda em aula que me gerou preocupação enquanto
\ac{ages} II. Reportei a situação à orientadora Paola após a retrospectiva:

\begin{quotation}
\itshape
``Boa noite. Estive conversando com o pessoal sobre o banco de dados e a
situação ficou um pouco confusa para todos. Havíamos montado o banco juntos em
call, com antecedência, justamente para evitar problemas. No entanto, após algumas
sugestões de mudança, atualizamos o \ac{uml} e o banco conforme o proposto, mas as
novas alterações não foram aceitas inicialmente. Isso acabou gerando uma nova task de
integração, pela qual o Tarcísio ficou responsável.

Na última quarta-feira (08/04), alguns campos foram alterados em conversa com o Gui
e o Felipe. Já nesta segunda-feira, o Felipe nos informou que utilizou a base do
Tarcísio, aplicou modificações e gerou uma nova versão do banco em uma branch ---
porém, não há registros de commits ou pushes para acompanharmos. O Tarcísio ficou
frustrado, sentindo que seu esforço foi em vão, e o restante do grupo ficou com a
sensação de que não pôde contribuir efetivamente com essa parte do projeto.

Falei com o Vini, que conseguiu contato com o Felipe, que prometeu realizar o commit
até amanhã (16/04). Caso haja algum imprevisto, nos reuniremos neste final de semana
para resolver. Por fim, gostaria de pontuar que o Heitor foi o único que não participou
de nada relacionado ao banco --- tive até que chamá-lo na sala de aula. O resultado da
tela de login que ele fez ficou muito bom, mas no banco não ajudou em nada.''
\end{quotation}

A Paola retornou imediatamente e informou que a questão do banco era decorrente de
uma combinação feita com o Felipe (\ac{ages} III), que havia assumido a responsabilidade
e não havia entregado, resultando em um banco não funcional. Ela também esclareceu
que a reprovação do nosso banco havia sido um erro que passou despercebido.

\textbf{Aprendizado:} não vale a pena trabalhar em um banco que não seja versionado
e com interação entre todos os integrantes. O \textit{dbdiagram} foi um obstáculo,
pois toda interação dependia de uma reunião ou de contato direto com o Tarcísio.
Como \ac{ages} III no próximo semestre, caso os \ac{ages} II sugiram esse \textit{framework},
defenderei o uso de uma plataforma que permita acesso e edição concorrente por todos
os membros --- como o Miro ou o Figma.
